\documentstyle[% 


righttag% 


,11pt% 


,amssymb% 


,russian%               remove in Eng. 


,izvru%                 Set izven% in Eng. 


]{amsart} 


 


%       Math definitions 


 


\newcommand{\A}{\cal A} 


\newcommand{\col}{\operatorname{col}} 


\newcommand{\vraisup}{\operatornamewithlimits{vrai~sup}} 


\newcommand{\tts}[1]{} 


 


%\hoffset=-15mm%                  For Eng. 


  


\setcounter{page}{28} 


\god{1999} 


\nomer{4~(443)} %                        Remove in Eng. 


%\nomer{Vol.\,43, No.\,4}%               Set in Eng. 


%\str{\pageref{begin}--\pageref{end}}%  Set in Eng. 


\udk{517.956} 


 


\author{\it В.Ф.\,Волкодавов, В.Н.\,Захаров} 


% NB:   ^^ remove in Eng. 


\title{Экстремальные свойства решений одного 


уравнения гиперболического типа третьего 


порядка в трехмерном пространстве  


и их применение} 


 


\date{} 


% \pagestyle{myheadings}%         Set in Eng. 


\pagestyle{plain} %              Remove in Eng. 


\begin{document} 


% \thispagestyle{plain}%          Set in Eng. 


\maketitle 


\label{begin} 


 


 


Известно [1], [2], что экстремальные свойства решений дифференциальных 


уравнений позволяют доказать единственность решений краевых 


задач для этих уравнений. Такие свойства были доказаны для 


ряда уравнений второго порядка с двумя независимыми переменными 


[3], [4]. В данной работе рассматривается уравнение третьего 


порядка в трехмерном пространстве 


\begin{equation} 


\cal U_{xyz}=0. 


\end{equation} 


 


Введем обозначения 


\begin{gather}


\begin{align*}


G^+&=\{(x,y,z):0<x-y<z<h,\ h>0\},\quad 


G^-=\{(x,y,z):0<z<x-y<h\},


%\notag


\\ 


G_1&=\{(x,y,z):y=0,\ 0\le x\le z\le h\},\quad\qquad 


G_2=\{(x,y,z):x=h,\ 0\le h-y\le z\le h\},


%\notag


\\ 


G_3&=\{(x,y,z):z=x-y,\ 0\le y\le x\le h\},\quad 


G_4=\{(x,y,z):z=0,\ 0\le y\le x\le h\},


%\notag


\end{align*}


\\


% 


\begin{align*}


H^+&=\{(x,y,z):0<x<z<h,\ 0<x<y<+\infty\},


%\notag


\\


H^-&=\{(x,y,z):0<x<z<h,\ 0<y<x<h\},


%\notag


\end{align*}


\\


H_1=\{(x,y,z):y=x,\ 0\le x\le z\le h\},\quad 


H_2=\{(x,y,z):z=x,\ 0\le x\le\ y\le h\},\notag \\ 


H_3=\{(x,y,z):x=0,\ 0\le y\le+\infty,\ 0\le z\le h\},\notag \\


H_4=\{(x,y,z):z=h,\ 0\le x\le h,\ 0\le y<+\infty\}, \notag \\


\begin{align} 


\cal U(x,y,z)\big|_{G_1}&=f_1(x,z), \\ 


\cal U(x,y,z)\big|_{G_2}&=g_1(y,z), \\ 


\cal U(x,y,z)\big|_{G_3}&=\tau_1(x,y), \\ 


\cal U_y(x,y,z)\big|_{G_3}&=\omega(x,z), \\ 


\cal U(x,y,z)\big|_{G_4}&=\varphi_1(x,y), \\ 


\cal U(x,y,z)\big|_{H_1}&=\tau_2(x,z), 


\end{align} 


\\ 


\frac{\partial\cal U}{\partial\ov m}\Big|_{H_1}= 


\bigg(\frac{\partial\cal U}{\partial x}- 


\frac{\partial\cal U}{\partial y}\bigg)\Big|_{H_1}= 


\nu(x,z), 


\end{gather} 


где $\frac{\partial\cal U}{\partial\ov m}$ --- производная по нормам к 


плоскости $y=x$, $\ov m=(1,-1,0)$, 


\begin{align} 


\cal U(x,y,z)\big|_{H_2}&=g_2(x,y), \\ 


\cal U(x,y,z)\big|_{H_3}&=\varphi_2(y,z), \\ 


\cal U(x,y,z)\big|_{H_4}&=f_2(x,y). 


\end{align} 


Для доказательства экстремальных свойств решений уравнения (1)  


воспользуемся единственным решением некоторых краевых задач в 


областях $G^-$, $G^+$, $H^-$ и $H^+$. Доказательства существования и 


единственности 


решений этих краевых задач здесь приводить не будем. Области,  


краевые условия и единственные решения представим в виде таблицы. 


\medskip 


 


{\hfil 


\begin{tabular}{|c|c|c|c|c|} 


\hline 


\No & Задача & Область & Краевые & Единственное решение задачи \\ 


\ & \ & \ & условия & \ \\


\hline 


1. & $D$\,I & $\overset{\phantom{.}}{G}{}^+$ & (2)--(4)


& $\cal U_+(x,y,z)=\tau_1(x,y)+ 


f_1(x,z)-f_1(x,x-y)+$ \\ 


&&&& $+g_1(y,z)-g_1(0,z)-g_1(y,x-y)+g_1(0,x-y)$ \\ 


2. & $B$ & $G^-$ & (4)--(6) & $\cal U_-(x,y,z)=\tau_1(x,x-z)+ 


\varphi_1(x,y)-$ \\ 


&&&& $-\varphi(x,x-y)+\int\limits^{x-z}_y\varphi_{1t} 


(t+z,t)dt-\int\limits^{x-z}_y\omega(t+z,t)dt$\\ 


3. & $C-G$ & $H^-$ & (7)--(9) & $\cal U_-(x,y,z)= 


\frac{1}{2}[\tau(x,z)+\tau_2(y,z)]+g(x,y)-$ \\ 


&&&& $-\frac{1}{2}\int\limits^y_x[\nu(t,z)-\nu(t,x)]dt$ \\ 


4. & $D$\,II & $H^+$ & (7), (10), (11) & $\cal U_+(x,y,z)= 


\tau_2(x,z)+f(x,y)-f(x,h)+$ \\ 


&&&& $+\varphi(y,z)-\varphi(x,z)-\varphi(x,h)- 


\varphi(y,h)$ \\ 


\hline 


\end{tabular} 


\hfil} 


 


\begin{lem} Если решение $\cal U_+(x,y,z)$ уравнения $(1)$ в области 


$G^+$ 


таково, что функция $\cal U(x,y,x-y)=\tau_1(x,y)$


достигает наибольшего  


положительного $($наименьшего отрицательного$)$ значения во внутренней 


точке $M_0$ множества $G_3$ и при этом $\cal U_+(x,0,z)$ на множестве 


$G_1$ и $\cal U_+(h,y,z)$ на множестве $G_2$ тождественно равны нулю, 


то 


$$ 


\frac{\partial\cal U_+(M_0)}{\partial\ov n}=0,\quad 


\frac{\partial^2\cal U_+(M_0)}{\partial\ov n^2}<0 


\ \; \bigg( 


\frac{\partial^2\cal U_+(M_0)}{\partial\ov n^2}>0\bigg), 


$$ 


где $\frac{\partial\cal U}{\partial\ov n}=-\cal U_x+\cal U_y+\cal U_z$, 


$\frac{\partial^2\cal U}{\partial\ov n^2}=\cal U_{xx}+\cal U_{yy}+\cal 


U_{zz}-2\cal U_{xy}-2\cal U_{xz}+2\cal U_{yz}$, $\ov n=(-1,1,1)$. 


\end{lem} 


 


\begin{pf} Пусть для решения задачи $D$\,I выполнены условия  


леммы. Тогда\linebreak $\cal U_+(x,y,z)=\tau_1(x,y)$. Вычисления дают 


$$ 


\frac{\partial\cal U_+}{\partial\ov n}\Big|_{z=x-y}= 


\tau_{1y}(x,y)-\tau_{1x}(x,y). 


$$ 


Tак как функция $\tau_1(x,y)$ в точке $M_0$ имеет экстремум, то 


$\frac{\partial\cal U_+(M_0)}{\partial\ov n}=0$. 


Нетрудно получить 


\begin{equation} 


\frac{\partial^2\cal U_+(M_0)}{\partial\ov n^2}= 


\tau_{1xx}(M_0)+\tau_{1yy}(M_0)-2\tau_{1xy}(M_0). 


\end{equation} 


 


Введем обозначения $\tau_{1xx}(M_0)=a$, $\tau_{1yy}(M_0)=b$, 


$\tau_{1xy}(M_0)=c$. Пусть в  


точке $M_0$ достигается положительный максимум. Тогда по достаточному 


условию существования экстремума функции двух переменных  


$a<0$, $ab-c^2>0$. Очевидно, $b<0$ и $\sqrt{ab}>|c|$. С учетом этих 


неравенств равенство (12) примет вид 


$$ 


\frac{\partial^2\cal U_+(M_0)}{\partial\ov n^2}=- 


(\sqrt{-a}-\sqrt{-b}\,)^2-2(\sqrt{ab}+c). 


$$ 


Отсюда следует 


$\frac{\partial^2\cal U(M_0)}{\partial\ov n_2}<0$. 


Аналогично убеждаемся, что 


$\frac{\partial^2\cal U_+(M_0)}{\partial\ov n^2}>0$, если в точке 


$M_0$ 


достигается отрицательный минимум.  


\end{pf} 


 


Уравнение (1) рассмотрим на множестве $G=G^-\cup G^+$. 


\medskip 


 


{\bf Задача $\Gamma$.} {\it 


Найти функцию $\cal U(x,y,z)$ со следующими свойствами$:$ 


\begin{itemize} 


\item[1)] $\cal U(x,y,z)\in\Bbb C(\ov G);$ 


\item[2)] $\cal U(x,y,z)$ --- решение уравнения $(1)$ на множестве $G$, 


удовлетворяющее 


условиям $(2)$, $(3)$, $(6)$ и 


условиям сопряжения 


\end{itemize}


}


\begin{align} 


\frac{\partial\cal U}{\partial\ov n}\Big|_{z=x-y+0}&= 


\frac{\partial\cal U}{\partial y}\Big|_{z=x-y-0}, \\ 


\frac{\partial^2\cal U}{\partial\ov n^2}\Big|_{z=x-y+0}&= 


\bigg[\,\frac{\partial^2\cal U}{\partial y\partial z}- 


\frac{\partial\omega(z+y,y)}{\partial(z+y)}\bigg] 


\Big|_{z=x-y-0}. 


\end{align} 


 


\begin{thm} Если существует решение задачи $\Gamma$, то оно 


единственно. 


\end{thm} 


 


\begin{pf} Достаточно показать,  


что решение задачи $\Gamma$ 


с однородными краевыми условиями тождественно 


равно нулю. Пусть функции $f_1(x,y)$, $g_1(x,y)$ и $\varphi_1(x,y)$ 


тождественно 


равны нулю. Из условия~1) постановки задачи $\Gamma$ следует, что  


функция $\cal U(x,y,x-y)=\tau_1(x,y)$ непрерывна на множестве $G_3$. 


По теореме 


Вейерштрасса на этом множестве она достигает своих наибольшего 


и наименьшего значений. Пусть она принимает свое наибольшее  


положительное значение в точке $M_0(x_0,y_0,x_0-y_0)$ множества $G_3$. 


Так как функция $\tau(x,y)$ на границе множества $G_3$ обращается в 


нуль, то  


точка $M_0$ внутренняя. Но тогда по лемме~1 


$$ 


\frac{\partial^2\cal U(M_0)}{\partial\ov n^2}<0. 


$$ 


Из решения задачи $B$ получаем 


$$ 


\bigg[\, 


\frac{\partial^2\cal U}{\partial y\partial z}- 


\frac{\partial\omega(z_y,y)}{\partial(z+y)}\bigg] 


\Big|_{z=x-y-0}\equiv0, 


$$ 


а это противоречит условию сопряжения (14). Значит, внутри множества 


$G_3$ 


функция $\tau(x,y)$ не может достигать своего наибольшего  


положительного значения. Следовательно, на множестве $G_3$ функция 


$\tau(x,y)=C$ --- const. Так как на границе $G_3$ функция $\tau(x,y)=0$, 


то  


$\tau(x,y)\equiv0$ на $G_3$. Но тогда из единственности решения задачи 


$D$\,I 


следует, что $\cal U(x,y,z)\equiv0$ в $\ov G^+$. Из условия~1) 


постановки 


задачи $\Gamma$ 


и условия сопряжения (13) получаем $\cal U(x,y,z)\equiv0$ на множестве 


$\ov G\,^-$.  


Отсюда $\cal U(x,y,z)\equiv0$ в $G$. 


\end{pf} 


 


\begin{lem} Если решение $\cal U_-(x,y,z)$ уравнения $(1)$ в области 


$H^-$ 


таково, что функция $\cal U_-(x,y,z)=\tau_2(x,z)$ достигает своего 


наибольшего 


положительного $($наименьшего от\-ри\-ца\-тель\-но\-го$)$ значения во 


внутренней точке $M_0(x_0,y_0,z_0)$ множества $H_1$ и при этом $\cal 


U_-(x,y,x)$ 


и $\frac{\partial\cal U_-(x,y,z)}{\partial\ov m}$ тождественно равны 


нулю на множествах $H_1$ и $H_2$ 


соответственно, то 


$$ 


\frac{\partial^2\cal U_-(M_0)}{\partial\ov m^2}<0\quad 


\bigg(\,\frac{\partial^2\cal U_-(M_0)}{\partial\ov m^2}>0 


\bigg).


$$ 


\end{lem} 


 


\begin{pf} При выполнении условий леммы 2 решение задачи  


$C-G$ примет вид 


$$ 


\cal U_-(x,y,z)=\frac{1}{2}[\tau(x,z)+\tau(y,z)]. 


$$ 


Отсюда 


$$


\frac{\partial^2\cal U(x,y,z)}{\partial\ov m^2}=\frac{1}{2} 


\bigg[\,\frac{\partial^2\tau(x,z)}{\partial x^2}+ 


\frac{\partial^2(y,z)}{\partial y^2}\bigg]. 


$$


Пусть в точке $M_0(x_0,y_0,z_0)$ достигается положительный максимум. 


Тогда 


$$ 


\frac{\partial^2\cal U_-(M_0)}{\partial\ov m^2}= 


\tau_{xz}(x_0,z_0)<0 


$$ 


в силу достаточного признака существования экстремума функции  


двух переменных. Аналогично, если в точке $M_0$ достигается 


отрицательный 


минимум, то 


$\frac{\partial^2\cal U_-(M_0)}{\partial\ov m^2}>0$. 


\end{pf} 


 


\begin{lem} Если решение $\cal U_+(x,y,z)$ уравнения $(1)$ в области $H^+$ 


таково, что функция $\cal U_+(x,x,z)$ достигает своего наибольшего 


положительного 


$($наименьшего отрицательного$)$ значения во внутренней 


точке множества $H_1$ и при этом $\cal U_+(0,y,z)$ и 


$\cal U_+(x,y,h)$ 


тождественно 


равны нулю на множествах $H_3$ и $H_4$ соответственно, то 


$$ 


\frac{\partial\cal U_+(M_0)}{\partial\ov m}=0,\quad 


\frac{\partial^2\cal U_+(M_0)}{\partial\ov m^2}<0 


\ \; \bigg(\, 


\frac{\partial^2\cal U_+(M_0)}{\partial\ov m^2}>0\bigg). 


$$ 


\end{lem} 


 


Доказательство этого утверждения проводится аналогично доказательству 


леммы 2 с использованием решения задачи $D$\,II. 


\medskip 


 


{\bf Задача $D$.} {\it  


Пусть $H=H^-\cup H^+$. Найти функцию $\cal U(x,y,z)$ 


со следующими свойствами$:$ 


\begin{itemize} 


\item[1)] $\cal U(x,y,z)\in\Bbb C(\ov H);$ 


\item[2)] $\cal U(x,y,z)$ --- решение уравнения $(1)$ на множестве $H$, 


удовлетворяющее условиям $(9)$--$(11)$ и 


условиям сопряжения 


\end{itemize}


}


\begin{align*} 


\frac{\partial\cal U_-}{\partial\ov m}\Big|_{y-x=-0}&= 


\frac{\partial\cal U_+}{\partial\ov m}\Big|_{y-x=+0}, \\ 


\frac{\partial^2\cal U_-}{\partial\ov m^2}\Big|_{y-x=-0}&= 


-\frac{\partial^2\cal U_-}{\partial\ov m^2} 


\Big|_{y-x=+0}. 


\end{align*} 


 


\begin{thm} Если существует решение задачи $D$, то оно 


единственно. 


\end{thm} 


 


Доказательство этой теоремы аналогично доказательству теоремы 1. 


 


\begin{thebibliography}{9} 


 


\bibitem{1} 


Смирнов М.М. {\em Уравнения смешанного типа}. -- М.: Высш. школа,  


1985. -- 304 с. 


\bibitem{2} 


Волкодавов В.Ф., Николаев Н.Я. {\em Краевые задачи для уравнения  


Эйлера-Пуассона-Дарбу}. -- Куйбышев.: Изд-во пед. ин-та, 1984. -- 80 с. 


\bibitem{3} 


Волкодавов В.Ф. {\em Об одной формулировке локального  


экстремума} // Волжск. матем. сб. -- 1970, вып.\,11. -- С.\,42--54. 


\bibitem{4} 


Волкодавов В.Ф., Невойструев Л.М. {\em О принципе локального 


экстремума для уравнения Эйлера-Пуассона-Дарбу и его применение} // 


Волжск. матем. сб. -- 1966, вып.\,5.


 


\end{thebibliography} 


 


\vskip 2cm 


 


 


\postup{\it Самарский государственный}{$\qquad$\it Поступили} 


\postup{\it педагогический университет}{{\it первый вариант} 20.11.1995} 


\postup{}{{\it окончательный вариант} 05.01.1999} 


\label{end} 


\end{document} 


